\section*{Guía de lectura y niveles de evidencia}
\addcontentsline{toc}{section}{Guía de lectura y niveles de evidencia}

La memoria separa lo demostrado, lo observado y lo pendiente. Cada afirmación
fuerte remite a una prueba o a datos regenerables; su alcance queda restringido
por el modelo, la información y la planta declarados. Esta convención evita
transferir a hardware una conclusión obtenida en simulación o confundir un
equilibrio estratégico con factibilidad mecánica.

\begin{table}[!htbp]
  \centering
  \caption*{Clave de estados empleada en la memoria.}
  \small
  \renewcommand{\arraystretch}{1.14}
  \begin{tabularx}{\textwidth}{@{}p{2.55cm}p{4.10cm}X@{}}
    \toprule
    Estado & Soporte exigido & Lectura permitida \\
    \midrule
    \textsc{Soportada} & prueba completa o datos y análisis reproducibles & conclusión válida solo en el dominio declarado \\
    \textsc{Parcial} & evidencia incompleta o indirecta & resultado acotado; no sostiene la conclusión fuerte \\
    \textsc{Pendiente} & falta verificación decisiva & hipótesis abierta, sin uso concluyente \\
    \textsc{Refutada} & contraejemplo o gate incumplido & resultado negativo conservado como evidencia \\
    \bottomrule
  \end{tabularx}
\end{table}

Un \emph{teorema} incluye supuestos y demostración; una \emph{proposición}
declara un alcance más estrecho; una \emph{conjetura} no se presenta como
hecho; una \emph{observación empírica} informa configuración, semillas y
estadística. Las magnitudes físicas usan unidades SI. Las cifras del
Capítulo~6 proceden de macros generadas, y las tablas distinguen explícitamente
los oráculos globales de los mecanismos distribuidos.

\begin{center}
  \begin{tikzpicture}[>=Latex, node distance=1.0cm, every node/.style={font=\small}]
    \node[draw, rounded corners, inner sep=5pt] (claim) {afirmación};
    \node[draw, rounded corners, inner sep=5pt, right=of claim] (assume) {supuestos};
    \node[draw, rounded corners, inner sep=5pt, right=of assume] (support) {prueba o datos};
    \node[draw, rounded corners, inner sep=5pt, right=of support] (scope) {alcance};
    \draw[->] (claim) -- (assume);
    \draw[->] (assume) -- (support);
    \draw[->] (support) -- (scope);
  \end{tikzpicture}
\end{center}

Esta cadena es el contrato de lectura del documento: si falta un eslabón, el
estado se degrada y la conclusión se limita de forma explícita.
